% ---
% id: receiver/lem:fock-monotone
% kind: lem
% title: "Order monotonicity of the Fock lift"
% summary: >-
%   The Fock lift preserves the Loewner order on positive semidefinite
%   operators. A larger detection operator A therefore gives a smaller lifted
%   no-click operator, the lift of 1 − A.
% prominence: standard
% status: accepted
% uses: [lem:fock-functor]
% assumes: []
% refs: []
% provenance: newly-derived
% lean:
% ---

\begin{lemma}\label{lem:fock-monotone}
If $X\succeq Y\succeq0$ on $\mathcal{K}_0$ then $\Fock(X)\succeq\Fock(Y)$.
Consequently, if $0\preceq A\preceq A'\preceq\id$ then
$\Fock(\id-A)\succeq\Fock(\id-A')$.
\end{lemma}

\begin{proof}
Telescoping,
\begin{equation}\label{eq:telescope}
  X^{\otimes n}-Y^{\otimes n}
  =\sum_{j=0}^{n-1}X^{\otimes j}\otimes(X-Y)\otimes Y^{\otimes(n-1-j)} .
\end{equation}
Each summand is a tensor product of positive semidefinite operators and is
therefore positive semidefinite, so $X^{\otimes n}\succeq Y^{\otimes n}$ on
$\mathcal{K}_0^{\otimes n}$. Both $X^{\otimes n}$ and $Y^{\otimes n}$ preserve
$\Symn{n}$ by Lemma~\ref{lem:fock-functor}\,(i), so for
$\ket{\psi}\in\Symn{n}$ the difference of the restrictions is
$\bra{\psi}(X^{\otimes n}-Y^{\otimes n})\ket{\psi}\ge0$: the restriction of a
positive semidefinite operator to an invariant subspace is positive
semidefinite. Summing over $n$ gives $\Fock(X)\succeq\Fock(Y)$. The second
statement is the first with $X=\id-A$, $Y=\id-A'$, both positive semidefinite
because $A,A'\preceq\id$.
\end{proof}
